\scanpage{79}
% archive foliation: 78
applique la prop. à $f'$, et on applique $=$ à $f''$.

\deleted{\unsure{Prop}} \uline{Corollaire 2}
Sous les conditions de la prop. pour $W$,
\deleted{$\mathfrak X$ \unsure{et}} supposons que
\deleted{$\mathfrak X\unsure{\times}\mathfrak Y$}
tte partie finie de $\mathfrak X$ soit contenue dans un ouvert affine.
Alors si $f:\mathfrak X\to\mathfrak Y$ est séparé, alors il en est
de $=$ de $\underline W^a(f)$ [sous réserve que
$\underline W^a(\mathfrak Y)$ existe dans le cas $\mathrm I'$].

On applique la prop. à $\mathfrak X\unsure{-|\to}\mathfrak X\times_{\mathfrak Y}\mathfrak X$.

\uline{Proposition} $W$, $\mathfrak Y$ comme dans la prop. précédente.
Soit $f:\mathfrak X\to\mathfrak Y$, on suppose
\deleted{$\unsure{\mathfrak X}$ \unsure{quasi}}
\inserted{dans le cas $\mathrm I'$}, pour \unsure{tout} ouvert affine
\inserted{$\mathcal U$} de $\mathfrak Y$,
\inserted{tte partie finie de} $f^{-1}(\mathcal U)$ est
\deleted{\unsure{réunion}} contenue dans un ouvert affine
(p. ex. $f$ quasi-projectif).
On suppose de plus $\mathfrak Y$ loc. noeth et \deleted{$\mathfrak X$}.
$f$ de type fini. \inserted{\deleted{\unsure{resp. affine}}}
Alors $\underline W^a(f)$ est de type fini
\deleted{\unsure{quasi-affine}}.
\margin{\unsure{inquiétude} Dans le cas ``Greenberg spécial'',
\unsure{l'énoncé est} \unsure{valable} \unsure{compte tenu}
\unsure{de la propriété} \unsure{projectif} $\Rightarrow$ quasi-projectif}
